\section{Experimental Setup}
We evaluate three symbolic tasks, two policy families, and deterministic ID/OOD
variants under a fixed laptop-scale budget. Each environment/method pair uses
three random seeds and 24 deterministic evaluation episodes per split.
ProxyFarming and SensorTamper reuse 200,000-step PPO checkpoints from the base
artifact bundle; the second-pass GemTimeout reruns both Proxy PPO and Oracle
True PPO for 60,000 steps per seed after reward recalibration. Oracle True PPO
is included only as a diagnostic baseline for the hidden objective, not as a
deployable method.

Monitor evaluation is stricter than the policy benchmark. Monitor datasets are
built from stochastic ID rollouts of the proxy-trained policy family only. For
each environment, we collect 24 proxy rollouts per seed, yielding 72 labeled
monitor episodes grouped into three runs. Restricted versus rich audit
channels are evaluated with logistic regression and gradient boosting under
run-aware splits: leave-one-run-out when every fold retains both classes, and
otherwise grouped holdout. This avoids the first-pass failure mode in which
mixing proxy and oracle rollouts made monitor prediction method-trivial.

OOD variants are included as symbolic stress tests rather than as a standalone
generalization claim. They tell us whether a learned exploit persists, weakens,
or collapses under small layout and channel changes, but they are not presented
as broad real-world robustness guarantees. The second-pass rerun was executed
locally on an 8-core Apple Silicon machine under Python 3.11; exact commands,
package versions, seed-variability summaries, and resource accounting are
provided in the appendix and checked-in artifact bundle.

In line with recent warnings about statistical instability in deep-RL
evaluation \citep{agarwal2021precipice}, we keep the inferential claims modest:
the paper reports seed standard deviations and cross-seed dominance margins,
but does not claim fine-grained significance beyond what three seeds can
support.

\begin{table*}[t]
\centering
\scriptsize
\resizebox{\textwidth}{!}{%
\begin{tabular}{l l l l l l}
\toprule
Environment & Proxy reward & True objective & Corruption channel & Monitor channels & OOD shift \\
\midrule
ProxyFarming & Beacon reward & Reach goal efficiently & Repeated beacon harvesting & Restricted trace; rich adds coarse audit state & Beacon move; layout; distractors \\
SensorTamper & Sensor-mediated reward & Deliver object to target & Tamper tile spoofs sensor & Restricted trace; rich adds coarse audit state & Tamper move; delay; path change \\
GemTimeout & Gem reward & Reach goal before timeout & Gem detours trade off with time & Restricted trace; rich adds coarse audit state & Density; budget; deceptive placement \\
\bottomrule
\end{tabular}
}
\caption{Compact summary of the three benchmark environments used in the main paper.}
\label{tab:env-summary-compact}
\end{table*}

